\Needspace{18\baselineskip}
\section{Exact search transport and reflection through recursive fillings}
\label{sec:expansion-search}
The preceding section constructs recursive fillings and actual specimen states.
We now prove a complementary operational result: context search survives those
fillings with its exact length, and a filled-graph search step at an original
vertex either comes from the context or touches an original shallow hole.
This is a concrete traversal theorem for the coordinate search relation.
It also yields executions and rootedness of actual states obtained by inserting
a focus after filling. No output-closure hypothesis is needed for these results.

The scope of this milestone is precise. The state construction here first fills
a focus-free context and then inserts a focus at a transported coordinate.
The earlier deep specimen construction expands an already focussed context.
Their compatibility, and the identification of the boundary predicate below
with the paper's entering/exiting specimen predicates, remain to be proved.
Consequently these results do not yet discharge the abstract
\texttt{context\_step} field or all of Lemma~39(2--3).

\Needspace{14\baselineskip}
\subsection{Transporting individual search rules}
Fix a graph $C:I\to O$ at depth $d$ and a family $H$ of replacements for its
outer edges. Write $E=C[H]$ for simultaneous expansion and
$j_H:V_C\to V_E$ for its injective host-vertex map.
\begin{definition}[Search-preserving family]\label{term:search-preserving}
A family preserves search when every original operation edge and every original
variable edge is replaced by a literal corolla copy. No restriction is imposed
here on replacements of other labels. In particular, this predicate alone does
not assert that newly available steps can be reflected back to the context.
\end{definition}
\lean{Graph/ExpansionSearch}{15}{18}{Linear.ExpansionSearch.PreservesSearch}
For a search coordinate $p=(m,v)$, define $j_Hp=(m,j_Hv)$: the mode is unchanged.
\lean{Graph/ExpansionSearch}{19}{22}{Linear.ExpansionSearch.position}
An operation view consists of its edge, operation symbol, singleton source,
ordered argument list, and ordered thunk list. Given a copy of that edge, its
expanded view uses the copy edge and maps its source, arguments, and thunks
through $j_H$. The copy-incidence equalities prove the required source and
target equations; mapping preserves the argument-list length.
\lean{Graph/ExpansionSearch}{23}{34}{Linear.ExpansionSearch.operation}
\begin{lemma}[Entry and return commute with expansion]
The expanded operation's entry coordinate is $j_H$ of the original entry
coordinate. Its return coordinate after argument $i$ is likewise $j_H$ of the
original return coordinate.
\end{lemma}
\begin{proof}
For entry, split the argument list into the empty and nonempty cases. The empty
case keeps the source and uses the operation's finish mode; the nonempty case
selects its first mapped argument. For return, lookup at $i+1$ commutes with
list mapping. A successful lookup selects the mapped argument; a failed lookup
keeps the mapped source with the same finish mode.
\end{proof}
\lean{Graph/ExpansionSearch}{35}{39}{Linear.ExpansionSearch.operation_enter}
\lean{Graph/ExpansionSearch}{40}{46}{Linear.ExpansionSearch.operation_resume}
\begin{theorem}[Forward search transport]
If $H$ preserves search and $p\longrightarrow_C q$, then
$j_Hp\longrightarrow_E j_Hq$.
\end{theorem}
\begin{proof}
Consider the three constructors grouping the five search rules. A variable
trigger uses its retained copy and the copied singleton source. Operation entry
uses the expanded operation view and the entry lemma. Operation return maps the
witnessed argument lookup and applies the return lemma. The argument order and
the exclusion of thunk targets are preserved in each case.
\end{proof}
\lean{Graph/ExpansionSearch}{47}{67}{Linear.ExpansionSearch.step}
If a context step exists, determinism on $E$ forces every successor of $j_Hp$
to equal its transported result. This conclusion does not require reflection
of arbitrary steps at a hole boundary.
\lean{Graph/ExpansionSearch}{68}{75}{Linear.ExpansionSearch.successor_unique}

For a recursive paired edge certificate, an original operation or variable
cannot be a filled-hole edge and cannot be a box. Its certificate must therefore
be retained, supplying copies on both sides. Applying this argument at each
outer edge proves that both filling families preserve search. Changes inside
box bodies are allowed throughout this argument.
\lean{Graph/ExpansionSearch}{76}{88}{Linear.DeepFillingEdge.search_copies}
\lean{Graph/ExpansionSearch}{89}{95}{Linear.deepFilling_preservesSearch}

\Needspace{14\baselineskip}
\subsection{Exact executions and successful search costs}
\begin{definition}[Coordinate search machine]\label{term:coordinate-machine}
The states are all mode--vertex pairs of $C$, and the transition relation is
coordinate search. A successful final coordinate has backtrack mode at an input
vertex. This machine may be defined on a context containing holes; it is not
asserted to be a closed UAM state. No edge targets an input, so a final coordinate
has no successor. Search determinism gives determinism of this machine.
\end{definition}
\lean{Graph/ExpansionSearch}{96}{106}{Linear.coordinateMachine}
\lean{Graph/ExpansionSearch}{107}{117}{Linear.coordinateMachine_deterministic}
\begin{theorem}[Exact execution transport]
Every $n$-step context execution $p\longrightarrow_C^n q$ yields
$j_Hp\longrightarrow_E^n j_Hq$ for a search-preserving family.
\end{theorem}
\begin{proof}
Induct on the execution. The empty execution maps to the empty execution.
For a nonempty execution, transport its first step and prepend it to the
inductively transported tail. No extra step is inserted and no step is deleted.
\end{proof}
\lean{Graph/ExpansionSearch}{118}{125}{Linear.ExpansionSearch.steps}
The host map retains the mode and whether a vertex is an input; hence finality
of $j_Hp$ is equivalent to finality of $p$. This is checked for each of the
internal, input, and output vertex constructors.
\lean{Graph/ExpansionSearch}{126}{133}{Linear.ExpansionSearch.final_iff}
\begin{theorem}[Uniqueness of the successful cost]
Suppose a context execution from $p$ finishes successfully after $n$ steps.
Every successful expanded execution from $j_Hp$ has length $n$ and ends at the
transported context result.
\end{theorem}
\begin{proof}
Transport the context run and its finality. Compare the supplied successful run
with this run in the deterministic coordinate machine. The previously proved
uniqueness theorem for successful executions gives equality both of lengths
and of endpoints. A separate exclusion of hole boundaries is unnecessary
because a complete successful context run has already been supplied.
\end{proof}
\lean{Graph/ExpansionSearch}{134}{148}{Linear.ExpansionSearch.final_run_unique}

\Needspace{14\baselineskip}
\subsection{Realizing executions as actual states}
A typed coordinate of a closed-interface context is sent through $j_H$ with its
term-typing proof: host expansion preserves vertex labels. The initial search
coordinate remains the initial coordinate at the unique input.
\lean{Graph/ExpansionSearch}{149}{154}{Linear.TypedPosition.expand}
\lean{Graph/ExpansionSearch}{155}{161}{Linear.TypedPosition.expand_initial}
\begin{theorem}[Execution realization]
On a focus-free, hole-free graph, an $n$-step execution between typed coordinates
induces an $n$-step search execution between the recursive isomorphism classes
of their actual focussed realizations.
\end{theorem}
\begin{proof}
Induct on the coordinate execution. In the empty case the two typed coordinates
are equal by equality of their underlying pairs. In the successor case, search
preserves term typing, so the intermediate coordinate can be realized as an
actual state. The established coordinate-to-state search theorem supplies the
first state step, and the induction hypothesis supplies the tail. Literal
state search embeds into recursive state search, which is the relation on the
classes used here.
\end{proof}
\lean{Graph/ExpansionSearch}{162}{182}{Linear.TypedPosition.realize_steps}
An execution beginning at the initial coordinate consequently proves rootedness
of the realized endpoint. Its first realized state is initial; the concrete
characterization of rootedness by a run from an initial class applies.
\lean{Graph/ExpansionSearch}{183}{197}{Linear.TypedPosition.realize_rooted}

\Needspace{14\baselineskip}
\subsection{Reading a step back from the expanded graph}
\begin{definition}[Hole boundary]\label{term:hole-boundary}
A context coordinate $(m,v)$ is at a hole boundary if some original shallow
edge has a hole label and either $m$ is search and $v$ is one of that edge's
sources, or $m$ is backtrack and $v$ is one of its targets. This is a predicate
on the outer context graph. A hole strictly inside a box does not by itself
satisfy this predicate.
\end{definition}
\lean{Graph/ExpansionSearch}{198}{203}{Linear.ExpansionSearch.HoleBoundary}
\begin{lemma}[Incident host vertices determine incident owners]
If an expanded edge owned by $e$ has source $j_Hv$, then $v$ is a source of $e$.
The analogous statement holds for targets.
\end{lemma}
\begin{proof}
Unpack membership in the mapped incidence list of the inserted edge. An internal
vertex of the replacement cannot equal a host vertex because their carriers
are disjoint. An output cannot occur in a source list, and an input cannot
occur in a target list. In the remaining boundary case, injectivity of $j_H$
identifies $v$ with the corresponding original source or target port.
\end{proof}
\lean{Graph/ExpansionSearch}{204}{212}{Linear.ExpansionSearch.source_owner_mem}
\lean{Graph/ExpansionSearch}{213}{221}{Linear.ExpansionSearch.target_owner_mem}
\begin{lemma}[Reading back an operation copy]
An operation view whose edge belongs to a retained corolla has an original
operation view with matching source and ordered arguments under $j_H$.
\end{lemma}
\begin{proof}
The copy has only one edge, so its label identifies the original label as the
same operation. The existence-of-view theorem constructs an original view.
Transport that view through the copy. It and the supplied expanded view name
the same edge; uniqueness of operation data therefore identifies their sources
and argument lists. The proof normalizes the edge equality explicitly, so the
dependent copy types are not silently identified.
\end{proof}
\lean{Graph/ExpansionSearch}{222}{244}{Linear.ExpansionSearch.operation_readback}
\begin{definition}[Search origins]\label{term:search-origins}
A family has search origins if every operation or variable edge in a replacement
either belongs to a retained corolla copy or is owned by an original hole edge.
This supplements forward search preservation with the origin information needed
for reflection.
\end{definition}
\lean{Graph/ExpansionSearch}{266}{273}{Linear.ExpansionSearch.SearchOrigins}
Every recursive edge certificate has this origin property. A retained certificate
supplies a copy; a filled certificate supplies the original hole label. In the
boxed case, the single replacement edge still has a box label, so it cannot be
an operation or variable. The same argument on the right side proves the origin
property for both families of every recursive paired filling.
\lean{Graph/ExpansionSearch}{245}{265}{Linear.DeepFillingEdge.search_origin}
\lean{Graph/ExpansionSearch}{386}{405}{Linear.deepFilling_searchOrigins}

\begin{theorem}[Reflection or a hole boundary]
Assume search preservation and search origins. If $j_Hp\longrightarrow_E q$,
then either there is a context step $p\longrightarrow_C r$ with $q=j_Hr$, or
$p$ is at a hole boundary.
\end{theorem}
\begin{proof}
Inspect the interacting variable or operation edge. If its owner is a hole,
source-incidence reflection in the search cases, or target-incidence reflection
in the return case, gives the required boundary witness with its mode.
Otherwise the owner is retained. For a variable trigger, the original variable
has a singleton source, and the reflected source membership identifies it with
the original coordinate. For operation entry, read back the operation view and
use injectivity of the host map to recover the original source. For return,
read back the argument list; injectivity of the mapped optional lookup recovers
the same argument index in the original list. These cases produce an original
search step. Forward transport and determinism identify its image with the
supplied successor $q$.
\end{proof}
\lean{Graph/ExpansionSearch}{274}{332}{Linear.ExpansionSearch.reflect_or_boundary}
As an immediate corollary, away from hole boundaries, search on expanded graphs
is exactly the transported context relation. The forward implication rules out
the boundary alternative; the reverse implication is forward search transport.
\lean{Graph/ExpansionSearch}{333}{346}{Linear.ExpansionSearch.step_iff}

\Needspace{14\baselineskip}
\subsection{Reflection of entire executions}
Let $R$ be a predicate on context coordinates such that every coordinate in $R$
avoids hole boundaries and every context step from $R$ stays in $R$.
\begin{theorem}[Execution reflection in an invariant region]
Every $n$-step expanded execution starting at $j_Hp$ with $R(p)$ is the image
of an $n$-step context execution ending at a coordinate still in $R$.
\end{theorem}
\begin{proof}
Induct on $n$. For zero steps use equality of the endpoints. Otherwise reflect
the first step using the boundary exclusion, obtain its original successor,
and use closure of $R$ to apply the induction hypothesis to the tail. The
result retains the exact length and equality of the expanded endpoint with
the transported context endpoint.
\end{proof}
\lean{Graph/ExpansionSearch}{347}{365}{Linear.ExpansionSearch.steps_reflect}
\begin{theorem}[Exact successful costs when no hole is shallow]
If $C$ has no outer hole-labelled edge, then an $n$-step successful expanded
execution from $j_Hp$ exists if and only if an $n$-step successful context
execution from $p$ exists. Holes inside boxes are permitted.
\end{theorem}
\begin{proof}
Use the entire coordinate space for $R$: a boundary witness would exhibit a
forbidden shallow hole, and closure is automatic. Reflect the run and then
reflect finality. Conversely transport the given context run and its finality.
This is an equivalence for search-only executions; it says nothing about
compute or substitution transitions that might subsequently open or copy boxes.
\end{proof}
\lean{Graph/ExpansionSearch}{366}{385}{Linear.ExpansionSearch.final_runs_iff}

\Needspace{14\baselineskip}
\subsection{Paired actual executions and rooted endpoints}
Fix a focus-free context $C:[\mathsf{term}]\to[]$, recursive paired edge
certificates for $H,K$, and hole-free expanded results. For a typed context
coordinate $p$, define its left filled state by realizing $j_Hp$ on $C[H]$.
Recursive filling supplies global focus freedom, and the stated hole-freedom
certificate supplies closedness. The right construction uses $K$ in the same
way. Uniqueness of hole names and allowed-context membership are not needed
for these search theorems, so they are not required by this construction.
\lean{Graph/ExpansionSearch}{406}{410}{Linear.ExpansionSearch.filledState}
\begin{theorem}[Paired exact traversal]
An $n$-step context execution between typed coordinates yields $n$-step
executions between the corresponding actual state classes on both sides.
If the context execution starts at the initial coordinate, both endpoint
states are rooted.
\end{theorem}
\begin{proof}
The recursive certificates supply search preservation separately for $H$ and
$K$. Transport the same context execution through each family, then apply
execution realization. For rootedness apply the initial-run characterization
on each expanded graph; the host map retains the initial input. This does not
assert rootedness preservation for arbitrary UAM rewrite transitions.
\end{proof}
\lean{Graph/ExpansionSearch}{411}{425}{Linear.ExpansionSearch.paired_steps}
\lean{Graph/ExpansionSearch}{426}{437}{Linear.ExpansionSearch.paired_rooted}
Two executions of the same length from the same recursive state class have
the same endpoint: induction on the length uses search determinism to equate
the first successors, then compares the tails. Therefore an independently
supplied actual-state execution, of the length of a given context trace and
starting at its left filled state, must finish at the constructed left endpoint.
The result is about equality of state classes, not equality of chosen carriers.
\lean{Graph/ExpansionSearch}{438}{451}{Linear.ExpansionSearch.state_run_unique}
\lean{Graph/ExpansionSearch}{452}{461}{Linear.ExpansionSearch.actual_run_endpoint}

\Needspace{14\baselineskip}
\subsection{A deep example and the necessary boundary restriction}
Reuse the three-edge graph of Section~\ref{sec:deep-specimens}: a nullary
operation produces a thunk, a box contains the only hole, and a final operation
consumes the box's term output. This time start from its focus-free outer graph.
The filling retains both operations and replaces the box by the box containing
the previously constructed filled body. The edge certificates are retained at
the two operations and recursive at the box. Hole freedom follows from the
closed filled body and the absence of holes in the operation corollas.
\lean{Graph/ExpansionSearchExamples}{9}{16}{Linear.ExpansionSearchExamples.fillings}
\lean{Graph/ExpansionSearchExamples}{17}{26}{Linear.ExpansionSearchExamples.edges}
\lean{Graph/ExpansionSearchExamples}{27}{41}{Linear.ExpansionSearchExamples.closed}
The maker's operation view has no term arguments and one thunk target.
Let the returned coordinate be backtrack at the input. With passive operation
kinds, the nullary rule gives a one-step execution from the initial coordinate
to this returned coordinate; the thunk is not searched.
\lean{Graph/ExpansionSearchExamples}{42}{52}{Linear.ExpansionSearchExamples.maker}
\lean{Graph/ExpansionSearchExamples}{53}{54}{Linear.ExpansionSearchExamples.returned}
\lean{Graph/ExpansionSearchExamples}{55}{59}{Linear.ExpansionSearchExamples.context_run}
The paired-rootedness theorem gives a rooted actual state with the deep-filled
box still present. Its mode and input position prove finality directly. The
successful-cost uniqueness theorem shows that every successful coordinate run
from the expanded initial position has exactly one step and this endpoint.
\lean{Graph/ExpansionSearchExamples}{60}{63}{Linear.ExpansionSearchExamples.deep_return_rooted}
\lean{Graph/ExpansionSearchExamples}{64}{68}{Linear.ExpansionSearchExamples.deep_return_final}
\lean{Graph/ExpansionSearchExamples}{69}{77}{Linear.ExpansionSearchExamples.deep_return_cost}

For the negative example, take the one-hole corolla used as that box body.
Before filling, no search constructor applies: the sole label is neither a
variable nor an operation. After filling with the variable corolla, the inserted
variable triggers rewrite mode at the same host input. There is a hole-boundary
witness there, namely the original hole and its input source. There cannot be a
context step whose transport is this trigger, because the context has no search
step at all. Thus the boundary exception in reflection is essential, even for
a single linear hole. The example concerns triggering substitution; it does
not implement the substitution transition itself.
\lean{Graph/ExpansionSearchExamples}{78}{85}{Linear.ExpansionSearchExamples.hole_no_search}
\lean{Graph/ExpansionSearchExamples}{86}{92}{Linear.ExpansionSearchExamples.filled_hole_triggers}
\lean{Graph/ExpansionSearchExamples}{93}{96}{Linear.ExpansionSearchExamples.entering_hole}
\lean{Graph/ExpansionSearchExamples}{97}{103}{Linear.ExpansionSearchExamples.reflection_needs_boundary}
